\cleardoublepage

\chapter{外文翻译}
\noindent\textbf{论文题目：}\parbox[t]{0.86\linewidth}{追踪证据：为检索增强生成构建知识支撑的推理链}


\sectionnonum{摘要}
检索增强生成（Retrieval-Augmented Generation, RAG）为问答（Question Answering, QA）提供了一种有效范式：先用检索器召回相关文档，再用生成式阅读器基于文档生成答案。然而，现有检索器并不完美，经常在Top-$k$召回集合中混入与问题无关或“表面相关但实际上误导”的段落，从而向阅读器引入噪声（noise）并降低最终性能。该问题在多跳问题（multi-hop questions）中尤为突出：多跳问题需要多步推理（reasoning）把分散在多个文档中的证据片段组合起来，而噪声文档会显著破坏证据整合，削弱RAG模型的多跳推理能力。

为提升RAG在多跳场景下识别与整合支撑证据（supporting evidence）的能力，本文提出TRACE（consTructing knowledge-grounded ReAsoning Chains to identify and integrate supporting Evidence）。TRACE通过构造“以知识为支撑的推理链”（knowledge-grounded reasoning chains）来替代“直接拼接全部检索文档”的做法：先将检索文档转换为知识图谱（Knowledge Graph, KG），即由知识三元组（knowledge triples）组成的集合，三元组形式为$\langle$头实体（head entity），关系（relation），尾实体/属性值（tail entity/value）$\rangle$。知识三元组粒度更细、表达更凝练，每个三元组通常只承载一条事实性信息，因此相比句子或段落能更好地隔离无关信息、减少“长上下文迷失”（lost-in-the-middle）对证据识别的干扰。随后，TRACE利用自回归推理链构造器（Autoregressive Reasoning Chain Constructor）在问题条件下逐步选择三元组，形成逻辑连贯的推理链，用以连接多跳所需的关键事实。最后，TRACE要么直接把推理链作为上下文输入阅读器生成答案（TRACE-Triple），要么用推理链定位并筛选更相关的原始文档，再输入阅读器生成答案（TRACE-Doc）。

在三个多跳QA数据集上的零样本（zero-shot）实验表明：与直接使用全部检索文档作为上下文相比，TRACE在精确匹配（Exact Match, EM）上平均可提升最高14.03个百分点，并且推理链通常比完整文档更短，却足以支撑正确回答。这些结果说明，在RAG设定下，与其让阅读器在大量段落中“自己找证据”，不如先将证据结构化并构造更凝练、可组合的推理链，往往就能显著增强多跳推理表现。

\section{引言}
RAG模型在QA任务上取得了显著进展，其典型结构是“检索器—阅读器”（retriever-reader）架构：检索器负责从大规模语料（常见为Wikipedia）中召回与问题相关的一组文档，阅读器（通常由LLM实例化）在这些文档条件下生成答案。在零样本设定下，常见做法是把召回文档直接前置拼接到问题之前，再交给LLM生成答案。该流程虽然简单，但存在一个关键隐患：当召回集合中包含不相关内容时，阅读器会被噪声误导，产生错误推断或在证据不足时产生“高置信胡编”。已有研究显示，即使只混入少量无关上下文，LLM也可能被明显分散注意力并显著降低表现。

多跳问题进一步放大了噪声文档的破坏性。多跳问题往往需要至少2步推理：第一步从问题中涉及的某个实体或概念出发，找到中间实体或中间事实；第二步再基于该中间实体定位最终答案。若检索集合中混入大量无关文档，阅读器不仅要完成推理，还要在长上下文中进行证据筛选与对齐，推理链容易被打断，导致“能检索到信息但无法可靠组合”的失败模式。因此，提升多跳RAG能力的关键不只在于“检索更多”，而在于“识别并整合真正支撑答案的证据”。

本文提出TRACE以解决上述痛点。其核心思想是把“证据识别与证据整合”作为一个显式、可控的中间步骤：先将文档转换为知识三元组并构建局部KG，再从KG中构造推理链，把多跳所需的事实以逻辑连贯的三元组序列形式呈现给阅读器。这样的中间结构具有两点优势：（1）粒度更细、信息更凝练，有助于减少与问题无关的叙述性内容；（2）链式结构天然鼓励“证据连接”，使多跳推理更像沿着链条推进，而不是在段落集合中盲搜。最终，TRACE将推理链作为更高信噪比的上下文，显著提升RAG模型的多跳推理能力。

\section{问题定义}
本文聚焦多跳问题。记多跳问题为$q$，答案为$a$。检索器为每个问题返回$N$篇文档，记为$\mathcal{D}_q=\{d_1,d_2,\dots,d_N\}$。每篇文档通常由标题（title）与正文（text）组成。任务目标是在给定$q$与$\mathcal{D}_q$的情况下，生成正确答案$a$。

\section{方法}
\subsection{任务定义与总体思路}
TRACE整体流程可概括为三步：知识图谱生成、推理链构造、答案生成。给定$q$与$\mathcal{D}_q$，TRACE先生成局部KG$\mathcal{G}_q$，再基于$\mathcal{G}_q$构造推理链$z$，最后在$z$条件下生成答案$a$。可将其写为（符号与原文一致，仅作说明）：
\[
p(a\mid q,\mathcal{D}_q)\sim p(a\mid q,z,\mathcal{D}_q)\cdot p(z\mid q,\mathcal{G}_q)\cdot p(\mathcal{G}_q\mid \mathcal{D}_q),
\]
其中$p(\mathcal{G}_q\mid \mathcal{D}_q)$对应KG生成器，$p(z\mid q,\mathcal{G}_q)$对应推理链构造器，$p(a\mid q,z,\mathcal{D}_q)$对应阅读器。直观上，TRACE把“用文档回答”分解为“把文档变成可组合的事实单元 + 用事实单元构造证据链 + 在证据链上生成答案”。

\subsection{知识图谱生成（Knowledge Graph Generator）}
TRACE使用KG生成器把检索文档转换为知识三元组集合。其动机是：句子或段落通常包含多条事实，且夹杂大量与当前问题无关的描述；而知识三元组更细粒度、更“单事实”，更便于在后续步骤中筛选支撑证据。例如一句话“Albert Einstein (14 March 1879-18 April 1955) was a German-born theoretical physicist.”同时包含出生日期、死亡日期、国籍、职业等多条信息；对应的三元组“$\langle$Albert Einstein; date of birth; 14 March 1879$\rangle$”能把“出生日期”这一事实从句子中解耦出来，使后续证据选择更不易被无关信息干扰。

实现上，TRACE采用提示学习（in-context learning）方式，让LLM充当KG生成器，从每篇文档中抽取三元组。由于把所有检索文档拼接到一个长输入会加重长上下文的“中间遗忘”（lost-in-the-middle）问题，TRACE选择对每篇文档独立抽取三元组，然后通过“跨文档共享实体”在图层面形成跨文档联系。由于本文主要面向Wikipedia文档（包含标题与正文），TRACE将标题视作核心实体，并引导KG生成器主要抽取“标题实体—正文实体/属性”之间的关系三元组。这一设计降低了“任意实体两两建边”的难度，同时在经验上能取得较稳健的三元组质量。

\subsection{推理链构造（Autoregressive Reasoning Chain Constructor）}
在得到$\mathcal{G}_q$后，TRACE的目标是从图中选出能“逻辑连接”并共同支撑答案的三元组序列，即推理链（reasoning chain / reasoning path）。原文将推理链建模为自回归过程：设最大链长为$L$，第$i$步选择第$i$个三元组$z_i$，并根据此前选择的前缀$z_{<i}$与候选集合$\hat{\mathcal{G}}_i$决定下一步：
\[
p(z\mid q,\mathcal{G}_q)=\prod_{i=1}^{L}p(z_i\mid q,z_{<i},\hat{\mathcal{G}}_i),\quad
\hat{\mathcal{G}}_i=f(q,z_{<i},\mathcal{G}_q),
\]
其中$f(\cdot)$表示在全图中筛选候选三元组的“候选生成/打分”过程。TRACE将构造器拆成两部分：三元组排序器（triple ranker）与三元组选择器（triple selector）。

\textbf{三元组排序器（Triple Ranker）}负责从全图中挑出与当前状态更相关的一小批候选三元组，减少选择器面对的搜索空间。TRACE用双编码器（bi-encoder）实现排序器：将“问题$q$ + 已选三元组$z_{<i}$”拼接为查询，编码为向量$\mathbf{q}_i$；将每条候选三元组$t_j$编码为向量$\mathbf{t}_j$，并用内积$\mathbf{q}_i^\top\mathbf{t}_j$衡量相关性，取Top-$K$作为候选集$\hat{\mathcal{G}}_i$。原文形式为：
\[
\mathbf{q}_i=\text{Enc}(q,z_{<i}),\quad \mathbf{t}_j=\text{Enc}(t_j),\ \forall t_j\in\mathcal{G}_q,
\]
并据此取Top-$K$。

\textbf{三元组选择器（Triple Selector）}负责从$\hat{\mathcal{G}}_i$中选择下一条三元组，使得新三元组与已有链条构成连贯推理。TRACE用LLM充当选择器，并把选择任务表述为“多项选择题”（multiple-choice）：把每条候选三元组格式化为选项（例如B、C、D等），并加入一个特殊选项A表示“无需更多三元组”（no need for additional triples），用于自适应终止推理链。选择器被要求只输出一个选项字母，从而降低LLM在选择阶段产生三元组幻觉的风险（因为候选三元组全部来自KG）。此外，TRACE利用选择器对各选项token的logit定义分布，并结合束搜索（beam search）生成多条推理链：在每步扩展时保留概率最高的$b$个分支，最终得到$R$条推理链。原文指出，生成多条推理链有助于补全单条链可能缺失的关键信息，并在后续答案生成或文档筛选阶段提供更丰富、更稳健的证据集合。

\textbf{自适应终止（Adaptive Chain Termination）}是TRACE的重要设计：不同问题所需证据步数不同，固定链长容易引入冗余三元组。TRACE通过选项A让模型自行判定“链条已足够回答问题”并提前停止，从而在覆盖与噪声之间取得更优折中。实验显示，自适应链长显著优于固定链长设置。

\subsection{以推理链作为生成证据}
TRACE给出两种答案生成方式：
\begin{itemize}
    \item \textbf{TRACE-Triple：}直接把推理链三元组序列作为上下文输入阅读器，让阅读器基于这些三元组生成答案。三元组天然具备顺序（由构造器按步选择得到），因此直接按链条顺序拼接通常能得到较好的上下文组织。
    \item \textbf{TRACE-Doc：}使用推理链中的三元组反向定位其来源文档，从$\mathcal{D}_q$中筛选出更相关的文档子集，再把筛选后的文档作为上下文输入阅读器生成答案。由于多条三元组可能来自同一文档，TRACE-Doc使用“多数投票”（majority voting）方式给文档排序：每条三元组为其来源文档投票，按票数排序，票数为0的文档被过滤。
\end{itemize}
这两种方式分别对应“用三元组压缩上下文”与“用三元组过滤文档”。实验中，两者都明显优于直接使用全部文档的标准RAG。

\section{实验}
\subsection{数据集与评测指标}
本文在三个多跳QA数据集上评测：HotpotQA、\allowbreak 2Wiki\-MultiHopQA、\allowbreak MuSiQue。数据集通常需要2--4跳推理。每个问题对应一组已检索文档（来自Wikipedia），其中HotpotQA与2Wiki\-MultiHopQA每题给10篇文档，MuSiQue每题给20篇文档。评测指标使用EM与F1（与数据集标准一致）。由于测试集不可公开，实验遵循常见做法在各数据集的开发集上报告结果，并从训练集中抽取少量样本作为调参开发集。

\subsection{对比方法与设置}
基线方法分为四类：（1）朴素基线：不使用文档（仅问题），以及使用全部文档（即标准RAG）；（2）双编码器（bi-encoder）对文档排序后取Top-$M$作为上下文，如Contriever、DRAGON+、E5、E5-Mistral；（3）交叉编码器（cross-encoder）对文档重排后取Top-$M$，如monoT5、BGE、RankLLaMA；（4）链式思维检索（Chain-of-Thought, CoT）类基线IRCoT：交替生成CoT句子与检索，从而获取更相关的文档集。为公平比较，TRACE与各基线使用相同阅读器生成答案；基线中Top-$M$由开发集选择最优。

实现细节方面，本文使用LLaMA3-8B-Instruct同时作为KG生成器与三元组选择器；使用e5-mistral-7b-instruct作为三元组排序器；阅读器默认使用LLaMA3-8B-Instruct并采用零样本提示生成答案，同时也报告在Mistral与Gemma等不同阅读器上的泛化结果。超参数方面，最大链长$L$默认设为4；推理链条数$R$与束宽$b$默认设为5；候选三元组数$K$从$\{15,20,25,30\}$中选择在开发集表现最好的设置。

\subsection{实现流程与关键细节补充}
为便于理解TRACE在工程上如何落地，本节按“KG生成$\rightarrow$候选三元组筛选$\rightarrow$多项选择式三元组选择$\rightarrow$多链束搜索$\rightarrow$答案生成”的顺序补充关键细节（对应原文第3节与附录算法）。

\textbf{（1）KG生成的输入拆分与离线化思路。}\par
论文强调KG生成与问题无关：先对文档做三元组抽取，再在问答阶段使用这些三元组构造推理链。直接把$\mathcal{D}_q$所有文档拼成长输入会触发长上下文问题，导致模型忽略中间部分信息。TRACE因此对每篇文档独立抽取三元组，再通过共享实体在图层面形成跨文档连接。该策略也更利于工程落地：若文档来自固定语料库（如Wikipedia），可提前离线抽取三元组并缓存，在线推理时只需做推理链构造与答案生成。

\textbf{（2）候选三元组筛选（Top-$K$）为何必要。}从统计看，每个问题对应的KG往往包含数十到上百个实体与三元组，且图整体很稀疏，若直接让选择器在全图上做多步选择，会引入大量无关候选，既增加LLM推理负担，也提升误选概率。双编码器排序器将“问题+已选链条前缀”作为查询，在全图中取Top-$K$三元组作为候选集合$\hat{\mathcal{G}}_i$，从而把选择器的决策空间限制在更可能相关的区域。论文在消融中表明，移除排序器会显著降低EM/F1，说明Top-$K$筛选是TRACE能够稳定构造链条的关键环节。

\textbf{（3）多项选择式三元组选择与分布定义。}TRACE把选择器任务表述为多项选择题：将每条候选三元组渲染为一个选项（B、C、D…），并加入A表示“无需更多三元组”，用于自适应终止。选择器只输出一个字母，避免模型直接“编造三元组”。此外，论文并非只取argmax，而是使用每个选项token的logit来构造分布（概念上可理解为对这些logit做softmax），再与束搜索结合生成多条推理链。该设计有两个意义：其一，多链（$R$条）可以覆盖不同的推理路径，减少单链遗漏关键hop的风险；其二，分布信息让束搜索能够保留“次优但可能关键”的备选分支，从而在后续步骤补全证据链。

\textbf{（4）多链束搜索与链长自适应终止。}论文给出推理链构造伪代码：对每个问题生成$R$条链，每步将每条链扩展为$b$个候选，按累积概率选Top-$R$保留。复杂度近似为$\mathcal{O}(N+LR)$（忽略候选集内编码与选择器调用的常数项），其中$N$为文档数、$L$为最大链长、$R$为链条条数。链长方面，固定链长会引入冗余事实并增加噪声，因此TRACE加入“自适应终止”选项A：当选择器认为已有链条足以回答问题时输出A并停止扩展当前链。这使得平均链长在不同数据集上稳定在约3左右，而不是被最大链长强制拉长。

\textbf{（5）TRACE-Triple与TRACE-Doc的区别与排序问题。}TRACE-Triple直接用三元组链作为上下文，链条顺序天然提供了较好的输入组织。TRACE-Doc则要把链条中的三元组映射回其来源文档；当多个三元组来自同一文档时，如何排序这些文档会影响阅读器表现。论文采用多数投票：每条三元组给其来源文档投票，按票数降序排列，票数为0的文档直接过滤。该排序方式在实践中表现稳定，并显著降低“文档噪声率”（即无关文档占比）。

\section{结果与分析}
本文主要围绕五个问题报告结果并分析。

\textbf{（RQ1）TRACE整体效果如何？}在三个数据集上，TRACE-Triple与TRACE-Doc都取得最优或接近最优表现。相对“使用全部文档”的标准RAG，TRACE-Triple与TRACE-Doc在EM上平均提升分别可达14.03与13.46个百分点。相对强基线IRCoT，两者在EM上也有显著提升。一个值得注意的现象是，TRACE-Triple使用的上下文token数量显著少于直接使用文档或重排后的文档，但仍能达到更高的EM/F1，这支持了论文的核心结论：在RAG中，构造凝练推理链往往比把全部文档拼接给阅读器更有效。

从具体数值看，论文在三套数据集上报告了阅读器为LLaMA3时的整体结果。以“使用全部文档”的标准RAG为参照，HotpotQA上EM约为45左右，而TRACE-Triple与TRACE-Doc分别提升到50+与55左右；2WikiMultiHopQA上，标准RAG的EM在20多到30之间，而TRACE-Triple可提升到40+甚至45左右；MuSiQue上，标准RAG的EM通常在十几到二十出头，而TRACE-Triple与TRACE-Doc可以提升到30+。从上下文长度看，标准RAG往往需要上千token（因为需要拼接10或20篇文档），而TRACE-Triple的推理链上下文通常只有百级token；尽管上下文更短，性能却更高，说明“减少噪声并提高信息密度”比“提供更多原始段落”更关键。

论文还比较了不同阅读器（reader）下的泛化：将阅读器替换为Mistral或Gemma后，TRACE相对标准RAG的优势依旧成立，表明推理链是一种相对模型无关（reader-agnostic）的证据表示，能跨不同开源LLM稳定发挥作用。这一点对工程落地很重要：链条构造器与阅读器可以解耦，系统可以根据成本与效果选择不同阅读器，而不必担心推理链只对某个特定模型有效。

\textbf{（RQ2）把文档变成KG三元组的价值是什么？}论文构造了两个变体：用“句子”或用“文档”作为链条的基本单元，让构造器在句子/文档级别做链式选择。结果显示，这两种变体在所有数据集上都明显弱于“三元组链”（TRACE-Triple）。原因在于句子/文档仍包含多条信息与大量无关内容，难以像三元组那样实现“单事实粒度”的噪声隔离；三元组更利于在链条中精确对齐支撑证据，降低无关叙述对推理的干扰。

论文通过消融进一步说明：当把链条基本单位从三元组退化为句子时，性能会明显下降；当进一步退化为文档级链条时，性能下降更明显。该结论与直觉一致：句子虽然比文档更细，但仍可能包含多个事实与背景性描述，且句子之间的连接关系未必显式；文档更是包含大量与问题无关的冗余信息。三元组的优势在于“每条三元组只表达一个事实”，并且天然包含实体与关系，使链条连接更明确，从而更适合用来表达多跳推理所需的事实拼接。

\textbf{（RQ3）自回归推理链构造器是否必要？}论文将构造器替换为“从KG中直接检索Top-$T$三元组”作为上下文，发现性能明显下降。这表明多跳证据整合不仅是“挑出若干相关三元组”，更需要利用已选证据提示下一步证据选择：前缀$z_{<i}$为后续选择提供线索，使链条更可能闭合，避免遗漏关键hop。

具体来说，“Top-$T$三元组”方案相当于把问题当作单次检索，选出若干最相关事实，但这些事实之间不一定能形成连贯的证据链，也不一定覆盖多跳推理的所有中间实体。自回归构造则会在每一步把已选三元组作为上下文，促使下一步选择与链条相连的事实，从而更可能形成“实体可连接、逻辑可推进”的推理路径。论文结果显示，移除构造器会在三个数据集上带来显著降分，说明“链条构造”是TRACE提升多跳推理的核心机制之一。

\textbf{（RQ4）排序器与选择器分别贡献什么？}移除三元组排序器（让选择器在全KG上直接选）会显著降分，说明候选集合控制对于降低噪声与提高决策效率至关重要；移除三元组选择器（直接取排序器Top-$b$）同样会显著降分，说明LLM在“保证链条连贯性”方面提供了关键推理能力。论文也比较了不同排序器（E5、DRAGON+等）与不同选择器（LLaMA3、Mistral、Gemma等），发现选择与排序模型都会影响最终性能，而默认组合表现最好。

排序器的作用主要体现在“把无关候选挡在外面”。当移除排序器时，候选集合等于全KG，选择器面对的选项会混入大量与当前问题或链条前缀无关的三元组，导致错误选择概率上升。选择器的作用则体现在“链条连贯性与终止判断”：即便候选集合较相关，仍需要判断下一步应该选哪条事实才能与此前事实形成可解释的推理链，同时还需要在证据充分时输出A提前终止，避免冗余事实进入上下文。论文的消融表明，直接用排序器Top-$b$拼链（等价于不使用选择器推理）会显著降分，说明“仅靠相似度排序”不足以构造高质量多跳链条。

\textbf{（RQ5）最大链长$L$如何影响结果？}随着$L$增大，平均链长不会线性增长，而是逐渐趋缓，原因在于自适应终止策略会在“链条已足够”时提前停止。性能通常先升后降：当$L$过小，链条覆盖不足；当$L$过大，冗余或弱相关三元组被引入，反而混淆阅读器。该现象进一步强调了“凝练性”的重要性。

\textbf{补充：关键消融与统计现象。}论文进一步用消融实验把TRACE的收益拆解到“表示粒度、链条构造、候选控制、选择推理”四个维度，并报告了若干一致的现象。首先，当把三元组替换为句子或文档作为链条单元时，EM/F1在三个数据集上均出现明显下降，说明粒度粗化会把噪声重新引入上下文，削弱证据对齐能力。其次，当不使用链条构造器而直接取Top-$T$三元组时，性能同样显著下降，尤其在2WikiMultiHopQA与MuSiQue这类证据分散的数据集上更明显，说明链条前缀提供的“下一步线索”对于覆盖所有hop至关重要。再次，移除排序器或选择器都会严重退化：排序器负责把无关候选降到可控范围，选择器负责把候选组织成可解释链条并判断终止时机，两者缺一不可。

从统计角度，论文还报告了KG与推理链的结构特征：在HotpotQA与2Wiki\-MultiHopQA上，平均每题KG包含数十个实体与三元组，但图密度很低（高度稀疏），跨文档实体之间往往缺少显式关系；即便如此，构造出的推理链平均长度约为3左右，且用推理链筛选出的相关文档数平均在2--3篇之间，远少于原始提供的10或20篇文档。这也解释了TRACE-Doc为何能显著降低噪声：原始文档集合的无关文档比例很高，而推理链筛选后无关比例明显下降，阅读器因此更容易在更短上下文里完成多跳推理。

\textbf{补充：链条条数$R$、候选规模$K$与示例选择（demonstrations）。}论文在附加实验中讨论了若干超参数对性能的影响。其一，推理链条数$R$从1增大到若干值时，性能通常先提升后趋于稳定：多条链可以补充单条链遗漏的证据，但当关键证据已覆盖后，继续增加链条主要引入冗余。其二，候选规模$K$也呈现“先升后降”：$K$太小会遗漏关键三元组，$K$太大则把噪声带回候选集合，干扰选择器判断，因此需要在信息覆盖与候选噪声之间权衡。其三，论文强调示例选择对in-context learning很关键：移除示例（只给指令不提供demonstrations）会显著降分；使用固定示例也往往弱于“按相似度自适应检索示例”的策略，说明示例需要与当前问题类型匹配，才能更好地引导选择器构造连贯链条。

\textbf{案例分析：}论文给出多个HotpotQA例子说明链条可解释性。例如问题“Albert Einstein的父亲何时出生？”对应推理链可为$\langle$Albert Einstein; father; Hermann Einstein$\rangle$，$\langle$Hermann Einstein; date of birth; ...$\rangle$；再如“某球员在2012--13拜仁赛季作为后腰引入，他的生日是什么？”链条先定位“新引入球员=Javi Martínez”，再定位“Javi Martínez的出生日期”，最终得到答案。链条使证据路径更清晰，也便于后续人工核查。

论文还给出若干有代表性的“是/否”与“属性对齐”问题示例。例如判断两位作家是否同国籍，推理链分别抽取两位作家的nationality三元组，再由阅读器据此输出“no”；又如询问两位人物共同职业，链条分别抽取occupation三元组并对齐到同一职业类别（novelist），输出“novelist”。这些例子说明推理链不仅适用于实体桥接类多跳问题，也能支持“并列对齐”的多跳组合推理。更重要的是，推理链输出的三元组本身就提供了可审计证据：如果答案错误，开发者可以直接检查链条中的哪一步事实缺失或被误选，从而更容易定位问题来源（KG抽取错误、排序器召回偏差、选择器误判或阅读器生成错误）。

\section{相关工作}
与TRACE最相关的研究方向主要包括三类。第一类是多跳RAG与交错检索推理：例如IRCoT（Interleaving Retrieval with Chain-of-Thought）通过生成CoT句子来驱动检索并逐步补全证据链，但CoT文本可能出现幻觉或错误中间结论，从而引入新的噪声。TRACE的不同之处在于，它的推理链三元组来自检索文档并被显式约束在KG候选之内，降低了“中间推理文本胡编”对后续检索与回答的连锁污染。第二类是“RAG抗噪声/证据选择”方向：大量工作关注无关文档对LLM的干扰以及如何筛选证据；TRACE通过把证据单元细化到三元组并用链条构造整合证据，提供了一种更结构化的证据选择方案。第三类是结合知识图谱的RAG：既有工作往往依赖外部既存KG（如Wikidata）或把KG视作可检索语料的一部分；TRACE则从检索到的文档中即时构建局部KG，再用于当前问题的推理链构造，避免了对全局KG覆盖率与更新性的依赖，更适合“只拿到一批文档但需要做多跳推理”的RAG场景。

此外，部分工作通过训练或强化学习的方式提升RAG对噪声上下文的鲁棒性，例如在训练阶段引入“无关文档干扰”并让模型学会拒绝或过滤，但这往往需要额外标注或训练资源。TRACE的特点在于完全采用in-context learning流程，不需要对LLM或检索器进行额外训练即可获得显著收益，从而更符合“快速试验、快速迁移”的工程需求。与此同时，TRACE把推理过程显式落在三元组链条上，也与“可解释推理”（interpretable reasoning）的研究方向契合：系统不仅输出答案，还输出能被检查的证据路径，为后续的事实核验、错误分析与交互式调试提供了更清晰的接口。

\section{讨论与局限性}
论文列出三点主要局限：（1）本文以Wikipedia文档为主，KG生成器主要抽取“标题实体—正文实体/属性”的关系，而非所有实体对之间的关系，这简化了任务但限制了KG表达的丰富性；（2）由于缺少标注数据，论文无法直接定量评估KG三元组与推理链质量，只能通过最终QA性能侧面衡量，并辅以统计与定性分析；（3）推理链构造器需要访问选择器对选项token的logit以定义分布并进行束搜索，而当选择器是黑盒LLM时可能无法获得logit；在这种情况下可退化为每步只选一个三元组，从而只生成单条链。

此外，从工程落地角度，TRACE在标准RAG流程中新增“KG抽取 + 链条构造”步骤，带来额外时间与成本，但其优势是显著减少了阅读器需要处理的上下文长度（推理链往往只需百级token而非千级token），并能显著降低文档噪声率，从而在端到端系统中可能实现“用结构化预处理换取更稳定的推理质量”。未来可进一步探索更轻量的KG抽取、更多源文档的适配提示、以及更强的链条构造策略与验证机制。

进一步来看，上述局限也提示了几个值得继续推进的方向。第一，标题中心（title-centric）的三元组抽取虽然稳健，但可能遗漏“正文实体之间”的关键关系，尤其当多跳推理需要跨实体桥接时，若桥接关系只存在于正文实体对之间，标题中心抽取可能无法覆盖。因此，未来可以探索在保持抽取稳定性的同时，引入更通用的实体对关系抽取与消歧策略，例如先用标题中心抽取建立骨架，再对高置信实体对做补充关系抽取，从而在成本可控的前提下提升图的连通性。

第二，缺少三元组/链条的金标准标注使得“模块级评估”困难，容易出现“端到端指标提升但不清楚原因”的情况。后续可以考虑构造小规模但高质量的标注集，用于评估三元组事实性（是否被文档支撑）、三元组相关性（是否有助于回答问题）、链条连贯性（相邻三元组是否在实体层面可连接、是否形成闭合证据链）等维度，并进一步研究如何把这些维度纳入自动化评价。

第三，选择器logit依赖限制了对黑盒模型的适配。若只能得到最终选项而无法获得分布信息，则束搜索与多链生成会退化，可能影响覆盖性与鲁棒性。一个可行方向是用“多次采样 + 自一致性（self-consistency）”近似分布，或用独立的轻量判别器对候选选项打分，从而在不访问logit的情况下恢复一定程度的多链探索能力。与此同时，也可以研究把“是否终止（A选项）”独立建模为一个二分类决策，使终止策略更稳定、对模型接口更友好。

最后，从系统角度，TRACE虽然减少了阅读器上下文长度，但在在线场景仍要付出KG抽取与链条构造的额外成本。如何在端到端延迟、吞吐与质量之间做权衡，可能需要结合缓存、增量抽取、以及对KG生成器/选择器进行蒸馏或量化等工程手段。总体而言，TRACE展示了“先结构化、再推理”的有效性，但将其推广到更复杂、更真实的数据源仍需要在抽取可靠性、评估体系与工程接口上持续完善。

\section{结论}
本文提出TRACE，通过“文档$\rightarrow$KG三元组$\rightarrow$自回归推理链$\rightarrow$答案生成”的流程，显式构造以知识为支撑的推理链，用以识别与整合支撑证据，从而显著增强RAG的多跳推理能力。实验表明，推理链不仅能提高EM/F1，还能以更少上下文token达到更好效果，说明在RAG设定下，构造凝练、可组合的证据链往往比直接使用全部文档更充分、更高效。更广泛地看，TRACE强调把“证据表示与证据连接”作为RAG系统的核心控制点，为构建更可解释、更可验证、更鲁棒的检索增强推理系统提供了可行路径。

从研究意义上，TRACE也提示了RAG系统设计的一条通用原则：当任务需要复杂推理时，应当尽量把“证据筛选与证据组织”显式化，并用结构化中间表示承载推理状态，而不是把所有负担都交给最终的生成模型去“端到端隐式完成”。这种“先结构化、再推理、再生成”的分层范式，有望推广到更广泛的知识密集型应用中，例如科研文献问答、法律条款推理与企业知识库决策支持等场景。

\endinput

\endinput

\sectionnonum{摘要}
大型预训练语言模型能够在其参数中存储大量事实性知识，并在下游任务微调后取得很强的效果。然而，这类模型对知识的访问与精确操控能力仍然有限，因此在知识密集型任务上的表现往往不如专门面向知识检索与利用的体系结构；同时，如何为模型的输出提供可追溯的证据来源，以及如何随着世界变化而更新模型的知识，仍然是未解决的研究问题。尽管带有可微分访问机制的外部非参数化记忆在抽取式任务上已有探索，但其在通用文本生成场景中的系统性研究仍然不足。

本文提出一种通用的微调范式——检索增强生成（Retrieval-Augmented Generation, RAG），用于将预训练的参数化记忆与显式的非参数化记忆结合到同一生成模型中。我们构建的RAG模型以预训练的序列到序列（seq2seq）模型作为参数化记忆，以Wikipedia的稠密向量索引作为非参数化记忆，并通过预训练的神经检索器访问该索引。我们比较了两种RAG形式：一种在整个输出序列上共享同一组检索到的段落，另一种允许不同token使用不同段落。我们在多项知识密集型NLP任务上对RAG进行微调与评测，在三项开放域问答任务上达到新的最好结果，并超过仅靠参数化记忆的seq2seq模型与任务特化的“检索—抽取”系统。在知识密集型生成任务中，RAG生成的文本更加具体、多样且更符合事实。

\section{引言}
预训练神经语言模型已经被证明能够从数据中学习并内化大量知识，从而在不依赖外部记忆的情况下，充当一种参数化的隐式知识库。然而，这种“把知识压缩进参数”的方式也存在显著局限：模型难以在不再训练的情况下扩展或修正记忆；其预测过程难以直接解释；并且可能生成与事实不一致的“幻觉”内容。相比之下，将参数化记忆与非参数化（即基于检索）的外部记忆相结合的混合模型，能够使知识被直接编辑、扩展，并且检索到的证据可以被检查与解释，从而缓解上述问题。

已有工作如REALM与ORQA将可微检索器与掩码语言模型结合，在开放域抽取式问答上取得了进展，但其主要关注的是抽取式任务。本文将“混合记忆”的思想扩展到NLP中最常用也最通用的序列到序列模型，并提出检索增强生成（RAG）作为一套通用的微调方案：RAG以预训练的seq2seq Transformer作为参数化记忆，以Wikipedia的稠密向量索引作为非参数化记忆，并通过预训练的神经检索器（Dense Passage Retriever，DPR）从索引中检索潜在文档，再由生成器（BART）在输入与潜在文档条件下生成输出。我们将检索到的文档视作潜变量，在训练时对其进行边缘化，并分别给出“整段共享文档”（RAG-Sequence）与“逐token选择文档”（RAG-Token）两种形式。

与大量从零训练、为单一任务构建特定外部记忆结构的工作不同，本文研究的关键设定是：参数化与非参数化两个组件均为预训练，并预先加载了大量知识；同时，由于检索访问机制本身已预训练，模型在不依赖额外训练的情况下也具备一定的知识访问能力。实验结果表明，这种将检索与生成紧密结合的混合记忆模型，在知识密集型任务上具有显著优势：RAG在Natural Questions、WebQuestions与CuratedTrec上达到新的最好结果，并在MS-MARCO问答生成、Jeopardy问题生成与FEVER事实验证等任务上表现强劲；此外，通过替换非参数化索引，我们可以在不重新训练模型的情况下更新其世界知识。

\begin{figure}[ht]
\centering
\includegraphics[width=\linewidth]{rag_architecture.pdf}
\caption{方法概览：将预训练检索器（查询编码器+文档索引）与预训练seq2seq生成器（BART）端到端结合。对查询$x$使用最大内积检索（MIPS）找到Top-$K$文档$z_i$，将$z$视为潜变量，并对不同检索文档条件下的生成概率进行边缘化。}
\label{fig:rag_arch}
\end{figure}

\section{方法}
RAG模型使用输入序列$x$检索文本段落$z$，并在生成目标序列$y$时将这些段落作为额外上下文。模型由两部分组成：（1）检索器$p_{\eta}(z\mid x)$，返回给定输入$x$时段落$z$的（Top-$K$截断）分布；（2）生成器$p_{\theta}(y_i\mid x, z, y_{1:i-1})$，在输入$x$、检索段落$z$以及历史生成token的条件下生成当前token。为端到端训练，我们将检索到的段落视作潜变量，并通过Top-$K$近似进行边缘化。

\subsection{RAG-Sequence 与 RAG-Token}
RAG-Sequence假设整个输出序列由同一段落负责，因此将段落作为单一潜变量进行边缘化：
\begin{eqnarray*}
p_{\text{RAG-Sequence}}(y\mid x)
&\approx& \sum_{z\in \text{top-}k(p(\cdot\mid x))} p_{\eta}(z\mid x)\, p_{\theta}(y\mid x,z) \\
&=& \sum_{z\in \text{top-}k(p(\cdot\mid x))} p_{\eta}(z\mid x)\prod_i p_{\theta}(y_i\mid x,z,y_{1:i-1}).
\end{eqnarray*}
RAG-Token允许每个token使用不同段落，因此逐token边缘化潜变量：
\begin{eqnarray*}
p_{\text{RAG-Token}}(y\mid x)
&\approx& \prod_i \sum_{z\in \text{top-}k(p(\cdot\mid x))} p_{\eta}(z\mid x)\, p_{\theta}(y_i\mid x,z,y_{1:i-1}).
\end{eqnarray*}
当把分类标签视作长度为1的目标序列时，上述两种形式在序列分类任务上等价。

\subsection{检索器：DPR}
检索组件基于双编码器结构：文档编码器将段落$z$映射为稠密向量$\mathbf{d}(z)$，查询编码器将输入$x$映射为向量$\mathbf{q}(x)$，并以最大内积检索作为近似最近邻搜索：
\[
p_{\eta}(z\mid x)\propto \exp\left(\mathbf{d}(z)^{\top}\mathbf{q}(x)\right),\quad
\mathbf{d}(z)=\mathrm{BERT}_d(z),\quad \mathbf{q}(x)=\mathrm{BERT}_q(x).
\]
Top-$k$检索可以表述为最大内积搜索（MIPS）问题，并可用近似方法在亚线性时间内求解。本文使用预训练的DPR初始化检索器并构建Wikipedia段落索引，该索引即为RAG的非参数化记忆。

\subsection{生成器：BART}
生成器可由任意编码器—解码器结构实现。本文选用BART-large作为预训练seq2seq模型，并通过简单拼接将输入$x$与检索段落$z$合并作为编码器输入，使生成过程在外部证据条件下进行。BART的参数被视为RAG的参数化记忆。

\subsection{训练}
给定训练样本对$(x_j, y_j)$，本文最小化负对数边缘似然$\sum_j -\log p(y_j\mid x_j)$，并使用Adam进行随机梯度优化。由于在训练中更新文档编码器会导致索引需要频繁重建，成本较高，本文固定文档编码器与索引，仅微调查询编码器与BART生成器。训练过程中不需要显式的“应检索哪篇文档”的监督信号。

\subsection{解码}
RAG-Token可被视作标准自回归生成器，其转移概率为
\[
p'_{\theta}(y_i\mid x,y_{1:i-1})=\sum_{z\in \text{top-}k(p(\cdot\mid x))} p_{\eta}(z\mid x)\,p_{\theta}(y_i\mid x,z,y_{1:i-1}),
\]
因此可直接用常规beam search解码。RAG-Sequence的序列概率不分解为标准的逐token形式，本文采用“逐文档beam search + 边缘化”的方式，并给出较彻底的解码与更快的近似解码策略。

\section{实验}
本文在多类知识密集任务上评估RAG。非参数化知识源统一使用Wikipedia（2018年12月dump），将文章切分为不重叠的100词段落，总计约2100万段落，并用FAISS构建最大内积检索索引。训练时每个输入检索Top-$k$段落，$k$在$\{5,10\}$中选择，测试时依据开发集调参。

\subsection{开放域问答}
开放域问答将问题与答案视为输入输出文本对$(x,y)$，直接最小化答案的负对数似然，并以Exact Match作为指标。本文与两类范式对比：一类为抽取式问答系统（从检索段落中抽取答案span），另一类为闭卷问答（不检索，仅依赖参数化记忆生成答案）。实验覆盖Natural Questions、TriviaQA、WebQuestions与CuratedTrec等常用数据集。

\subsection{抽象式问答}
为评估知识密集场景下的自由形式生成能力，本文使用MS-MARCO NLG任务：数据提供问题与人工标注的整句答案。实验不使用数据集中提供的gold passages，而把任务当作开放域抽象式问答，以检验RAG在没有显式证据提示时的检索与生成能力。

\subsection{Jeopardy问题生成}
本文进一步构造更具挑战性的开放域问题生成任务：给定实体（Jeopardy答案）生成对应的Jeopardy问题。该任务要求生成精准且可核验的事实陈述，适合检验模型的知识性生成质量。评测使用Q-BLEU-1，并辅以人工评价，从事实性与特异性两个维度进行成对比较。

\subsection{事实验证}
FEVER任务要求判断一个自然语言断言在Wikipedia上是否被支持、被反驳或信息不足。本文把标签映射为长度为1的输出序列进行训练，并强调不使用检索证据的中间监督，以贴近真实应用中缺乏检索监督信号的情况。

\section{结果与分析}
在开放域问答任务中，RAG在多项数据集上取得新的最好结果，并且展示出兼具“开放书本”检索范式的性能与“闭卷”生成范式的灵活性。与使用重排序器与抽取式阅读器的系统相比，RAG表明在一定条件下无需额外的重排序与抽取模块也能获得很强表现。此外，即使正确答案未被任何检索段落包含，RAG仍可能生成正确答案，这体现了参数化与非参数化记忆的互补：检索提供线索与上下文，参数化记忆补全缺失信息。

\begin{table}[ht]
\centering
\caption{开放域问答测试集结果（EM/Acc.，数值与原论文一致）。TriviaQA包含两种评测设置：Open-Domain与TQA-Wiki。}
\label{tab:qa_scores_cn}
\small
\setlength{\tabcolsep}{0.4em}
\begin{tabular}{lccccc}
\toprule
模型 & NQ & TQA(Open) & TQA(Wiki) & WebQ & CTrec \\
\midrule
T5-11B & 34.5 & - & 50.1 & 37.4 & - \\
T5-11B+SSM & 36.6 & - & 60.5 & 44.7 & - \\
REALM & 40.4 & - & - & 40.7 & 46.8 \\
DPR & 41.5 & 57.9 & - & 41.1 & 50.6 \\
\midrule
RAG-Token & 44.1 & 55.2 & 66.1 & 45.5 & 50.0 \\
RAG-Sequence & 44.5 & 56.8 & 68.0 & 45.2 & 52.2 \\
\bottomrule
\end{tabular}
\end{table}

在MS-MARCO抽象式问答与Jeopardy问题生成中，RAG生成文本相较BART更具体、更符合事实，并表现出更高的多样性。实验还对检索机制做了消融：冻结检索器会降低多类任务表现，说明端到端学习检索有实质收益；同时，以BM25替代稠密检索在实体密集的FEVER上可能更占优，而在开放域问答等任务中稠密可微检索优势更明显，提示不同任务对检索信号存在差异。

\begin{table}[ht]
\centering
\caption{生成与分类任务结果（数值与原论文一致）。Jeopardy使用BLEU-1（B-1）与Q-BLEU-1（QB-1）；MS-MARCO给出Rouge-L与BLEU-1；FEVER给出标签分类准确率。}
\label{tab:gen_cls_scores_cn}
\small
\setlength{\tabcolsep}{0.35em}
\begin{tabular}{lcccccc}
\toprule
模型 & B-1 & QB-1 & R-L & B-1 & FEVER-3 & FEVER-2 \\
\midrule
SotA & - & - & 49.8* & 49.9* & 76.8 & 92.2* \\
BART & 15.1 & 19.7 & 38.2 & 41.6 & 64.0 & 81.1 \\
RAG-Token & 17.3 & 22.2 & 40.1 & 41.5 & 72.5 & 89.5 \\
RAG-Sequence & 14.7 & 21.4 & 40.8 & 44.2 & 72.5 & 89.5 \\
\bottomrule
\end{tabular}
\end{table}

\begin{figure}[ht]
\centering
\includegraphics[width=\linewidth]{rag_posterior_plot.pdf}
\caption{RAG-Token在Jeopardy问题生成中的文档后验示例：对输入“Hemingway”，模型在生成不同token时对不同检索文档的后验概率会发生变化，表明生成过程能够从多篇证据中组合信息。}
\label{fig:rag_posterior}
\end{figure}

\begin{figure}[ht]
\centering
\includegraphics[width=\linewidth]{rag_retrieval_plots.pdf}
\caption{检索文档数量对性能的影响：增加检索文档数通常会改善开放域问答结果，但在不同RAG变体与任务上存在收益递减或峰值现象。}
\label{fig:rag_retrieval_plots}
\end{figure}

本文还展示了非参数化索引“热替换”带来的知识更新能力：用不同时间点的Wikipedia构建索引，模型能随索引切换而改变对世界事实（如国家领导人）的回答，从而无需重新训练即可更新知识。

\section{相关工作}
过往研究已在多种任务中验证了检索对性能的提升，例如开放域问答、事实核查、长文本问答、对话与机器翻译等。本文的贡献在于将这些面向单任务的检索增强成功经验统一到一个可通用微调的生成式框架中，展示单一检索增强架构在多任务上的适用性与强性能。

此外，通用预训练模型（如BART、T5）在不使用检索的情况下也能覆盖多种生成与判别任务，但其知识更新与可解释性受限。本文通过引入可学习的检索模块，扩展了统一架构可覆盖的任务空间。关于学习型检索，已有工作利用搜索、强化学习或潜变量方法优化检索模块以服务特定任务；本文则强调同一检索增强架构可在多类任务上通过微调达到强性能。

\section{讨论}
本文提出的RAG模型展示了参数化与非参数化记忆结合的优势：在知识密集任务上取得强结果，在生成任务上减少幻觉并提升事实性与特异性；同时，索引热替换提供了无需再训练的知识更新路径。未来可以探索两部分组件的联合预训练，从而在更大规模与更复杂任务上进一步提升检索与生成的协同效果，并更深入理解两类记忆在生成过程中的交互机理。

\section{结论}
本文提出检索增强生成（RAG）作为一种通用微调范式，将预训练seq2seq生成器的参数化记忆与可检索的外部非参数化记忆有机结合，通过把检索段落视作潜变量并在Top-$k$近似下进行边缘化，实现“检索—生成”端到端协同。实验结果表明，RAG在多项开放域问答与知识密集生成任务中取得强性能，并通过替换索引实现无需再训练的知识更新。后续工作可从联合预训练、更强检索与排序组件、以及更精细的证据使用与可验证机制等方向进一步提升其实用性与可信性。

\section{更广泛影响}
RAG通过更强的事实性支撑降低幻觉，并提供可检查的证据来源，因此在医疗、教育、科研辅助等场景具有潜在积极价值。然而，外部知识源（如Wikipedia）本身可能包含错误与偏见；同时，生成式模型仍可能被滥用于制造误导性内容、垃圾信息或钓鱼文本。因此，在部署此类系统时需要引入审计、过滤与对抗滥用的机制，并以负责任的方式评估与缓解风险。

\section{致谢}
作者感谢评审的建设性意见，并感谢Hugging Face在开源实现方面的支持与帮助，同时也感谢相关讨论与资助项目对研究工作的支持。
